% GENERATED FILE. Edit canonical lessons, then run npm run generate:books.
% canonical-source-hash: fnv1a64:ec05c8516e00deeb
% canonical-lessons: MW-C79-ghanto, MW-C79-minat, MW-C79-aj, MW-C79-kal, MW-C79-abar

\chapter{Hour, Minute, Today, Yesterday, Now}
\label{ch:mw-hour-minute-today-yesterday-now}
\hlchaptermodality{79}

\begin{chapteropening}
\textbf{By the end of this chapter:} \emph{I can say an hour, a minute, today, yesterday and now.}

Complete the last lesson of chapter 79: I can say an hour, a minute, today, yesterday and now.
\end{chapteropening}

\section[\texorpdfstring{\mw{घंटो} (ghaṇṭo)}{घंटो (ghaṇṭo)}]{\mw{घंटो} (ghaṇṭo) — an hour}

\label{lesson:MW-C79-ghanto}



\begin{quote}
Before the new one: say the Marwadi for a month, then the Marwadi for a year.
\end{quote}

\subsection*{You'll want to know: \mw{घंटो}}
\textbf{\mw{घंटो}} — \emph{ghaṇṭo} — \textquotedblleft{}an hour\textquotedblright{}.

\begin{cousinweb}
\textbf{\mw{घंटो}} (\emph{ghaṇṭo}) is a tadbhava, a word inherited from Sanskrit and reshaped by sound change. It comes from \emph{ghaṇṭā}, \textquotedblleft{}bell\textquotedblright{}. CDIAL gives no Marwari form in this entry (CDIAL 4421); the match is through Hindi \emph{ghaṇṭā}. The bell that struck the hours gave its name to the hour.
\end{cousinweb}

\begin{grammarlens}[title={What you've built}]
One more word for this chapter.
\end{grammarlens}

\subsection*{Your turn}
\begin{itemize}
\raggedright
  \item \emph{Say it:} \emph{ghaṇṭo}
  \item \emph{Say it:} \emph{ghaṇṭo}, once more
  \item \emph{Say it:} say \emph{ghaṇṭo}
  \item \emph{Recall:} say the Marwadi for a month, then the Marwadi for a year, then say \emph{ghaṇṭo} again
\end{itemize}

\begin{tcolorbox}[breakable,colback=teal!4,colframe=teal!35!black,title={Before you move on}]
What does \mw{घंटो} mean? (\textquotedblleft{}An hour\textquotedblright{}.) Say it once more.
\end{tcolorbox}
\section[\texorpdfstring{\mw{मिनट} (minaṭ)}{मिनट (minaṭ)}]{\mw{मिनट} (minaṭ) — a minute}

\label{lesson:MW-C79-minat}



\begin{quote}
Before the new one: say the Marwadi for a year, then the Marwadi for an hour.
\end{quote}

\subsection*{You'll want to know: \mw{मिनट}}
\textbf{\mw{मिनट}} — \emph{minaṭ} — \textquotedblleft{}a minute\textquotedblright{}.

\begin{cousinweb}
\textbf{\mw{मिनट}} (\emph{minaṭ}) is a loan from English \emph{minute}.
\end{cousinweb}

\begin{grammarlens}[title={What you've built}]
One more word for this chapter.
\end{grammarlens}

\subsection*{Your turn}
\begin{itemize}
\raggedright
  \item \emph{Say it:} \emph{minaṭ}
  \item \emph{Say it:} \emph{minaṭ}, once more
  \item \emph{Say it:} say \emph{minaṭ}
  \item \emph{Recall:} say the Marwadi for a year, then the Marwadi for an hour, then say \emph{minaṭ} again
\end{itemize}

\begin{tcolorbox}[breakable,colback=teal!4,colframe=teal!35!black,title={Before you move on}]
What does \mw{मिनट} mean? (\textquotedblleft{}A minute\textquotedblright{}.) Say it once more.
\end{tcolorbox}
\section[\texorpdfstring{\mw{आज} (āj)}{आज (āj)}]{\mw{आज} (āj) — today}

\label{lesson:MW-C79-aj}



\begin{quote}
Before the new one: say the Marwadi for an hour, then the Marwadi for a minute.
\end{quote}

\subsection*{You'll want to know: \mw{आज}}
\textbf{\mw{आज}} — \emph{āj} — \textquotedblleft{}today\textquotedblright{}.

\begin{cousinweb}
\textbf{\mw{आज}} (\emph{āj}) is a tadbhava, a word inherited from Sanskrit and reshaped by sound change. It comes from \emph{adya}, \textquotedblleft{}today\textquotedblright{}. CDIAL lists Old Marwari \emph{āji} under it (CDIAL 242).
\end{cousinweb}

\begin{grammarlens}[title={What you've built}]
One more word for this chapter.
\end{grammarlens}

\subsection*{Your turn}
\begin{itemize}
\raggedright
  \item \emph{Say it:} \emph{āj}
  \item \emph{Say it:} \emph{āj}, once more
  \item \emph{Say it:} say \emph{āj}
  \item \emph{Recall:} say the Marwadi for an hour, then the Marwadi for a minute, then say \emph{āj} again
\end{itemize}

\begin{tcolorbox}[breakable,colback=teal!4,colframe=teal!35!black,title={Before you move on}]
What does \mw{आज} mean? (\textquotedblleft{}Today\textquotedblright{}.) Say it once more.
\end{tcolorbox}
\section[\texorpdfstring{\mw{काल} (kāl)}{काल (kāl)}]{\mw{काल} (kāl) — yesterday; tomorrow}

\label{lesson:MW-C79-kal}



\begin{quote}
Before the new one: say the Marwadi for a minute, then the Marwadi for today.
\end{quote}

\subsection*{You'll want to know: \mw{काल}}
\textbf{\mw{काल}} — \emph{kāl} — \textquotedblleft{}yesterday; tomorrow\textquotedblright{}.

\begin{cousinweb}
\textbf{\mw{काल}} (\emph{kāl}) is a tadbhava, a word inherited from Sanskrit and reshaped by sound change. It comes from \emph{kālya}, \textquotedblleft{}timely; at dawn\textquotedblright{}. CDIAL lists Old Marwari \emph{kālī} under it (CDIAL 3104). Sanskrit \emph{kalya} already meant \textquotedblleft{}at dawn, tomorrow\textquotedblright{}; the Old Marwari row means \textquotedblleft{}tomorrow\textquotedblright{}. It is built on \emph{kāla}, \textquotedblleft{}time\textquotedblright{}.
\end{cousinweb}

\begin{grammarlens}[title={What you've built}]
One more word for this chapter.
\end{grammarlens}

\subsection*{Your turn}
\begin{itemize}
\raggedright
  \item \emph{Say it:} \emph{kāl}
  \item \emph{Say it:} \emph{kāl}, once more
  \item \emph{Say it:} say \emph{kāl}
  \item \emph{Recall:} say the Marwadi for a minute, then the Marwadi for today, then say \emph{kāl} again
\end{itemize}

\begin{tcolorbox}[breakable,colback=teal!4,colframe=teal!35!black,title={Before you move on}]
What does \mw{काल} mean? (\textquotedblleft{}Yesterday; tomorrow\textquotedblright{}.) Say it once more.
\end{tcolorbox}
\section[\texorpdfstring{\mw{अबार} (abār)}{अबार (abār)}]{\mw{अबार} (abār) — now}

\label{lesson:MW-C79-abar}



\begin{quote}
Before the new one: say the Marwadi for today, then the Marwadi for yesterday; tomorrow.
\end{quote}

\subsection*{You'll want to know: \mw{अबार}}
\textbf{\mw{अबार}} — \emph{abār} — \textquotedblleft{}now\textquotedblright{}.

\begin{cousinweb}
CDIAL has no entry for \textbf{\mw{अबार}} (\emph{abār}), so its history is not settled. Its first part looks like Hindi \emph{ab}, \textquotedblleft{}now\textquotedblright{}, which CDIAL lists under Sanskrit \emph{evam eva}, \textquotedblleft{}exactly so\textquotedblright{}, crossed with the pronoun stem \emph{a-} (CDIAL 2528), where Old Marwari has \emph{hiva}, \textquotedblleft{}now\textquotedblright{}. The rest of \textbf{\mw{अबार}} (\emph{abār}) is not explained there.
\end{cousinweb}

\begin{grammarlens}[title={What you've built}]
One more word for this chapter.
\end{grammarlens}

\subsection*{Your turn}
\begin{itemize}
\raggedright
  \item \emph{Say it:} \emph{abār}
  \item \emph{Say it:} \emph{abār}, once more
  \item \emph{Say it:} say \emph{abār}
  \item \emph{Recall:} say the Marwadi for today, then the Marwadi for yesterday; tomorrow, then say \emph{abār} again
\end{itemize}

\begin{tcolorbox}[breakable,colback=teal!4,colframe=teal!35!black,title={Before you move on}]
What does \mw{अबार} mean? (\textquotedblleft{}Now\textquotedblright{}.) Say it once more.
\end{tcolorbox}
